\chapter{Heavy-Light Decomposition}

A balanced binary search tree answers range queries on an \emph{array} in $\Oh(\log n)$ time: a
segment tree or a Fenwick tree turns ``sum of a contiguous block'' into a logarithmic walk. Trees in
the graph sense are not arrays, though. A query like ``what is the maximum edge weight on the path
between vertices $u$ and $v$?'' or ``add $5$ to every vertex in the subtree rooted at $w$'' has no
obvious contiguous home in any array. The path from $u$ to $v$ can bend through the tree in a shape
that depends on $u$ and $v$, and there are $\Th(n^2)$ such paths.

Heavy-light decomposition (HLD) is the bridge. It carves the tree into a small family of vertical
\emph{paths}, lays each path down as a contiguous segment of a single array, and proves that any
root-to-node path -- and therefore any node-to-node path -- crosses only $\Oh(\log n)$ of these
segments. Once the tree is sliced this way, every off-the-shelf range structure you already own
(segment tree, Fenwick tree, sparse table) becomes a tree-path structure. Path queries and path
updates cost $\Oh(\log^2 n)$; subtree queries fall out almost for free.

This chapter builds the decomposition from the size-of-subtree intuition, proves the
$\Oh(\log n)$-chains bound rigorously, traces the whole machine by hand on a concrete $11$-vertex
tree, gives algorithm2e pseudocode for the build and for both kinds of query, and nails down the
complexity. We close with the subtree trick and the connection to Euler tours, link-cut trees, and
the LCA structures of the neighboring chapters.

\chapterline

\section{The problem: queries that follow tree paths}
\label{ch19:problem}

Fix a rooted tree $T$ on $n$ vertices, root $r$. Each vertex (or each edge) carries a value, and we
want to support operations of two flavors along \emph{paths} and \emph{subtrees}:

\begin{itemize}
  \item \textbf{Path query} $\textsc{Query}(u,v)$: aggregate the values on the vertices (or edges)
        lying on the unique simple path from $u$ to $v$ -- sum, max, min, gcd, anything associative.
  \item \textbf{Path update} $\textsc{Update}(u,v,\delta)$: add $\delta$ to (or assign $\delta$ on)
        every vertex on that path.
  \item \textbf{Subtree query / update}: the same, restricted to the subtree rooted at a given $w$.
  \item \textbf{Point operations}: read or write a single vertex; a degenerate path with $u=v$.
\end{itemize}

\begin{intuitionbox}
On an array, ``the path between positions $i$ and $j$'' is just the index range $[i,j]$, and a
segment tree handles it in $\Oh(\log n)$. On a tree the path between $u$ and $v$ is geometric: go up
from $u$, go up from $v$, meet at the lowest common ancestor $\ell=\mathrm{LCA}(u,v)$, and the path is
$u\!\to\!\ell$ concatenated with $\ell\!\to\!v$. We would love to relabel the vertices so that every
such climb is a contiguous index range. No single labeling can do that for \emph{every} path -- but
heavy-light decomposition does it with only $\Oh(\log n)$ ranges per path, which is just as good.
\end{intuitionbox}

The naive approaches set the bar. Re-walking the path vertex by vertex costs $\Oh(n)$ per query: fine
for one query, fatal for $10^5$ of them. Precomputing all pairs costs $\Th(n^2)$ memory. We want
preprocessing near-linear and each operation polylogarithmic, while supporting \emph{updates}, which
rules out static precomputation alone.

\begin{focusbox}
The whole chapter rests on one structural fact. Color each tree edge \emph{heavy} or \emph{light} by
a rule that depends only on subtree sizes. Then any root-to-leaf walk uses at most
$\floor{\log_2 n}$ light edges. Heavy edges chain together into vertical paths; light edges are the
rare jumps between chains. ``Few light edges'' is exactly ``few chains crossed,'' and ``few chains''
is exactly ``few segment-tree ranges.''
\end{focusbox}

\section{Heavy and light: the size rule}
\label{ch19:rule}

For a vertex $v$, let $\mathrm{sz}(v)$ be the number of vertices in the subtree rooted at $v$
(counting $v$ itself). For a non-leaf $v$, its \emph{heavy child} is a child $c$ maximizing
$\mathrm{sz}(c)$; ties are broken arbitrarily but consistently (say, the first such child in adjacency
order). The edge $(v,c)$ to the heavy child is the \emph{heavy edge}; every other child-edge of $v$
is \emph{light}.

\begin{definition}[Heavy edge, light edge]
\label{ch19:def-heavy}
Let $v$ have children $c_1,\dots,c_k$ in the rooted tree. The heavy child is
$\arg\max_i \mathrm{sz}(c_i)$ (lowest index on ties). The edge from $v$ to its heavy child is heavy;
the edges from $v$ to its other children are light. Every vertex (except the root, which has no parent
edge) sits at the bottom of exactly one parent edge, which is heavy or light accordingly.
\end{definition}

\begin{definition}[Heavy path / chain]
\label{ch19:def-chain}
A maximal connected set of vertices joined by heavy edges is a \emph{heavy path} (or \emph{chain}).
Because each vertex has at most one heavy edge going down (to its heavy child) and at most one coming
from above (it is its parent's heavy child or not), the heavy edges form a disjoint union of vertical
paths. Each chain has a \emph{top} (highest vertex, whose parent edge is light or absent) and runs
downward through a sequence of heavy children.
\end{definition}

\begin{intuitionbox}
Picture water poured at the root, always flowing into the fattest child. The streams it traces are
the heavy chains: long vertical runs down the bulkiest part of the tree. A light edge is a step into a
small side-branch -- and stepping into a small branch is something you cannot do too many times,
because each light step at least halves the subtree size beneath you. That halving is the entire
argument.
\end{intuitionbox}

\subsection{The key lemma}

\begin{lemma}[Light edges halve the subtree]
\label{ch19:lem-halve}
If $(v,c)$ is a light edge (so $c$ is a non-heavy child of $v$), then
$\mathrm{sz}(c) \le \mathrm{sz}(v)/2$.
\end{lemma}

\begin{whybox}
Let $h$ be the heavy child of $v$. By the choice of heavy child, $\mathrm{sz}(c)\le\mathrm{sz}(h)$
for every other child $c$. Now $v$'s subtree contains $v$, the subtree of $h$, the subtree of $c$, and
possibly more, all disjoint, so
\[
  \mathrm{sz}(v)\;\ge\; \mathrm{sz}(h) + \mathrm{sz}(c) \;\ge\; \mathrm{sz}(c) + \mathrm{sz}(c)
  \;=\; 2\,\mathrm{sz}(c).
\]
Hence $\mathrm{sz}(c)\le\mathrm{sz}(v)/2$. The only fact used is that the heavy child is at least as
big as $c$; everything else only makes $\mathrm{sz}(v)$ larger. \qed
\end{whybox}

\begin{theorem}[Few light edges on any vertical path]
\label{ch19:thm-loglight}
On the path from the root to any vertex $u$, the number of light edges is at most
$\floor{\log_2 n}$. Consequently the root-to-$u$ path meets at most $\floor{\log_2 n}+1$ distinct
heavy chains.
\end{theorem}

\begin{whybox}
Walk from the root down to $u$. Whenever we traverse a light edge $(v,c)$, Lemma~\ref{ch19:lem-halve}
says the subtree size at least halves: $\mathrm{sz}(c)\le\mathrm{sz}(v)/2$. Heavy edges never increase
the subtree size as we descend either. Start with $\mathrm{sz}(\text{root})=n$ and end with
$\mathrm{sz}(u)\ge 1$. If the path uses $k$ light edges, then
\[
  1 \;\le\; \mathrm{sz}(u) \;\le\; \frac{n}{2^{k}},
\]
because each of the $k$ light edges contributes a factor of at least $2$ to the denominator and heavy
edges only shrink things further. Rearranging, $2^k\le n$, so $k\le\log_2 n$, i.e.
$k\le\floor{\log_2 n}$ since $k$ is an integer.

For the chain count: a vertical root-to-$u$ path is partitioned into maximal heavy runs separated by
light edges. With $k$ light edges there are at most $k+1$ heavy runs, and each run lies inside a single
chain. So the path touches at most $k+1\le\floor{\log_2 n}+1$ chains. \qed
\end{whybox}

\begin{trapbox}
The bound is about \emph{light} edges, not chains directly, and it is a property of \emph{vertical}
root-descending paths. A general $u\!\to\!v$ path is two vertical paths glued at the LCA, so it meets
at most $2(\floor{\log_2 n}+1)=\Oh(\log n)$ chains -- still logarithmic. A common mistake is to claim
each chain is short. Chains can be as long as $n$ (a path graph is one giant chain). What is bounded is
how \emph{many different} chains a path visits, not their lengths.
\end{trapbox}

\section{Laying the chains into an array}
\label{ch19:layout}

Knowing the chains is only half the job. To reuse a segment tree we must assign each vertex a
\emph{position} $\mathrm{pos}(v)\in\{0,1,\dots,n-1\}$ so that:

\begin{enumerate}
  \item every chain occupies a \emph{contiguous} block of positions, with the chain's top at the
        smallest position and proceeding downward; and
  \item each subtree occupies a \emph{contiguous} block of positions (this gives subtree queries for
        free).
\end{enumerate}

Both are achieved by a single DFS that, at each vertex, recurses into the heavy child \emph{first},
then into the light children, assigning positions in visit order. Visiting the heavy child first keeps
a chain's vertices consecutive; visiting all of a vertex's subtree before backtracking keeps each
subtree consecutive. This is just an Euler-style preorder with a deliberate child ordering.

\begin{intuitionbox}
Think of two passes. The first DFS computes $\mathrm{sz}(\cdot)$ and, while it is at it, records each
vertex's heavy child. The second DFS does the relabeling: it stamps the current vertex with the next
free position, then dives heavy-first so an entire chain gets a run of consecutive stamps before any
light branch starts. Each vertex also remembers the \emph{top} of its chain, which is what lets a
query jump from chain to chain.
\end{intuitionbox}

\subsection{A tree to carry through the chapter}

\begin{figure}[H]
\centering
\begin{tikzpicture}[
  every node/.style={circle,draw=nodeblue,thick,fill=blue!5,minimum size=8mm,inner sep=0pt,font=\small},
  heavy/.style={very thick,structurepurple},
  light/.style={dashed,edgegray},
]
\node (1) at (0,0) {1};
\node (2) at (-3,-1.4) {2};
\node (3) at (2.2,-1.4) {3};
\node (4) at (-4.2,-2.8) {4};
\node (5) at (-1.8,-2.8) {5};
\node (6) at (1.0,-2.8) {6};
\node (7) at (3.4,-2.8) {7};
\node (8) at (-2.4,-4.2) {8};
\node (9) at (-1.2,-4.2) {9};
\node (10) at (0.4,-4.2) {10};
\node (11) at (-1.2,-5.6) {11};

\draw[heavy] (1) -- (2);
\draw[light] (1) -- (3);
\draw[light] (2) -- (4);
\draw[heavy] (2) -- (5);
\draw[light] (5) -- (8);
\draw[heavy] (5) -- (9);
\draw[light] (5) -- (10);
\draw[heavy] (9) -- (11);
\draw[heavy] (3) -- (6);
\draw[light] (3) -- (7);
\end{tikzpicture}
\caption{The running tree, rooted at $1$. Thick purple edges are heavy; dashed gray edges are light.
We will compute the sizes, chains, and array positions below.}
\label{ch19:fig-tree}
\end{figure}

The adjacency (children, in the order we will scan them) is:
$1\!:\{2,3\}$, $2\!:\{4,5\}$, $3\!:\{6,7\}$, $5\!:\{8,9,10\}$, $9\!:\{11\}$, and $4,6,7,8,10,11$ are
leaves.

\begin{labbox}
\textbf{Subtree sizes (post-order).} Leaves first: $\mathrm{sz}(4)=\mathrm{sz}(6)=\mathrm{sz}(7)=
\mathrm{sz}(8)=\mathrm{sz}(10)=\mathrm{sz}(11)=1$. Then
\[
  \mathrm{sz}(9)=1+\mathrm{sz}(11)=2,\quad
  \mathrm{sz}(5)=1+\mathrm{sz}(8)+\mathrm{sz}(9)+\mathrm{sz}(10)=1+1+2+1=5,
\]
\[
  \mathrm{sz}(2)=1+\mathrm{sz}(4)+\mathrm{sz}(5)=1+1+5=7,\qquad
  \mathrm{sz}(3)=1+\mathrm{sz}(6)+\mathrm{sz}(7)=3,
\]
\[
  \mathrm{sz}(1)=1+\mathrm{sz}(2)+\mathrm{sz}(3)=1+7+3=11 = n.\ \checkmark
\]

\textbf{Heavy children.} At $1$: children $2\,(7)$ vs $3\,(3)$ $\Rightarrow$ heavy child $2$. At $2$:
$4\,(1)$ vs $5\,(5)\Rightarrow 5$. At $3$: $6\,(1)$ vs $7\,(1)\Rightarrow 6$ (tie, lowest index). At
$5$: $8\,(1),9\,(2),10\,(1)\Rightarrow 9$. At $9$: only $11\Rightarrow 11$. This matches the thick
edges in Figure~\ref{ch19:fig-tree}.

\textbf{Chains.} Following heavy edges from each chain top:
\[
  \text{Chain A}: 1\!\to\!2\!\to\!5\!\to\!9\!\to\!11,\qquad
  \text{Chain B}: 3\!\to\!6,
\]
\[
  \text{singletons}: \{4\},\ \{7\},\ \{8\},\ \{10\}.
\]
Their tops are $1,3,4,7,8,10$ respectively (a vertex is a chain top iff its parent edge is light or it
is the root).
\end{labbox}

\subsection{Positions from the heavy-first DFS}

We now run the second DFS from the root, stamping positions $0,1,2,\dots$, always recursing into the
heavy child before the light ones. Whenever we enter a vertex through a light edge (or it is the root),
that vertex starts a new chain, so we set its $\mathrm{top}$ to itself; a heavy step inherits the
parent's $\mathrm{top}$.

\begin{labbox}
Order of first-visits, heavy-first at every vertex:
\[
  1,\;2,\;5,\;9,\;11,\;8,\;10,\;4,\;3,\;6,\;7.
\]
Reading the stamps off:
\begin{center}
\begin{tabular}{c|ccccccccccc}
vertex $v$        & 1 & 2 & 5 & 9 & 11 & 8 & 10 & 4 & 3 & 6 & 7\\\hline
$\mathrm{pos}(v)$ & 0 & 1 & 2 & 3 & 4  & 5 & 6  & 7 & 8 & 9 & 10\\
$\mathrm{top}(v)$ & 1 & 1 & 1 & 1 & 1  & 8 & 10 & 4 & 3 & 3 & 7\\
\end{tabular}
\end{center}
Chain A is exactly positions $[0,4]$ -- one contiguous block, top at $0$. Chain B is positions
$[8,9]$. The singleton chains land at $5,6,7,10$. Figure~\ref{ch19:fig-array} shows the layout.
Notice subtrees are contiguous too: subtree of $5$
is $\{5,9,11,8,10\}$ at positions $\{2,3,4,5,6\}=[2,6]$, and subtree of $2$ is positions $[1,7]$. We
will lean on that for subtree operations in Section~\ref{ch19:subtree}.
\end{labbox}

\begin{figure}[H]
\centering
\begin{tikzpicture}[
  cell/.style={draw=edgegray,minimum width=8mm,minimum height=8mm,inner sep=0pt,font=\small},
  chainA/.style={fill=structurepurple!18},
  chainB/.style={fill=quotientorange!22},
  single/.style={fill=blue!6},
]
\foreach \p/\vtx/\sty in {0/1/chainA,1/2/chainA,2/5/chainA,3/9/chainA,4/11/chainA,
                          5/8/single,6/10/single,7/4/single,8/3/chainB,9/6/chainB,10/7/single}{
  \node[cell,\sty] (c\p) at (\p*0.85,0) {\vtx};
  \node[font=\scriptsize\color{edgegray}] at (\p*0.85,-0.62) {\p};
}
% chain brackets
\draw[decorate,decoration={brace,amplitude=4pt},structurepurple,thick]
  (-0.42,0.5) -- (3.82,0.5) node[midway,above=3pt,font=\scriptsize]{chain A: $1,2,5,9,11$};
\draw[decorate,decoration={brace,amplitude=4pt},quotientorange,thick]
  (6.38,0.5) -- (8.07,0.5) node[midway,above=3pt,font=\scriptsize]{chain B: $3,6$};
\end{tikzpicture}
\caption{The position array. The top row is the vertex stored at each position; the bottom row is the
position index. Chain A (purple) occupies the contiguous block $[0,4]$, chain B (orange) occupies
$[8,9]$, and the four singleton chains $\{8\},\{10\},\{4\},\{7\}$ fall at positions $5,6,7,10$.}
\label{ch19:fig-array}
\end{figure}

\begin{whybox}
Why is a chain contiguous? When the DFS first reaches a chain top $t$, it stamps $t$, then immediately
recurses into $t$'s heavy child (the next chain vertex), stamps it, recurses into \emph{its} heavy
child, and so on -- no light branch is touched until the heavy descent of the chain has been fully
stamped, because heavy-first ordering puts the heavy subtree entirely before any light subtree.
Therefore the chain's vertices receive a run of consecutive numbers starting at $\mathrm{pos}(t)$. Why
is a subtree contiguous? Standard preorder: a vertex's stamp is followed by the stamps of its entire
subtree before the DFS backtracks, regardless of child order.
\end{whybox}

\section{Building the decomposition}
\label{ch19:build}

Two DFS passes, then we hand the position array to whatever range structure we like. The first pass
computes sizes, parents, depths, and heavy children; the second assigns positions and chain tops.

\begin{algorithm}[H]
\caption{\textsc{Decompose}: build heavy-light positions}
\label{ch19:alg-build}
\DontPrintSemicolon
\KwIn{Rooted tree $T$ on $n$ vertices, root $r$, adjacency lists.}
\KwOut{Arrays $\mathrm{par},\mathrm{dep},\mathrm{sz},\mathrm{heavy},\mathrm{top},\mathrm{pos}$.}
\BlankLine
\textbf{Function} $\textsc{DfsSize}(v,p)$\;
\quad $\mathrm{par}[v]\gets p$;\ \ $\mathrm{sz}[v]\gets 1$;\ \ $\mathrm{heavy}[v]\gets -1$;\ \ $\mathrm{best}\gets 0$\;
\quad \For{$c\in \mathrm{children}(v)$ \textnormal{(neighbors} $\ne p$\textnormal{)}}{
  \quad $\mathrm{dep}[c]\gets \mathrm{dep}[v]+1$;\ \ $\textsc{DfsSize}(c,v)$;\ \ $\mathrm{sz}[v]\mathrel{+}=\mathrm{sz}[c]$\;
  \quad \If{$\mathrm{sz}[c] > \mathrm{best}$}{$\mathrm{best}\gets \mathrm{sz}[c]$;\ \ $\mathrm{heavy}[v]\gets c$}
}
\BlankLine
\textbf{Function} $\textsc{DfsHld}(v,\ t)$\quad\textnormal{($t$ = current chain top)}\;
\quad $\mathrm{top}[v]\gets t$;\ \ $\mathrm{pos}[v]\gets \mathrm{timer}$;\ \ $\mathrm{timer}\mathrel{+}=1$\;
\quad \If{$\mathrm{heavy}[v]\ne -1$}{$\textsc{DfsHld}(\mathrm{heavy}[v],\ t)$ \quad\textnormal{// stay on the chain}}
\quad \For{$c\in \mathrm{children}(v)$ \textnormal{with} $c\ne \mathrm{heavy}[v]$}{
  \quad $\textsc{DfsHld}(c,\ c)$ \quad\textnormal{// light edge: $c$ starts a new chain}
}
\BlankLine
$\mathrm{dep}[r]\gets 0$;\ \ $\textsc{DfsSize}(r,-1)$;\ \ $\mathrm{timer}\gets 0$;\ \ $\textsc{DfsHld}(r,r)$\;
build a segment tree (or Fenwick tree) over an array of length $n$ indexed by $\mathrm{pos}$\;
\end{algorithm}

\begin{complexitybox}
$\textsc{DfsSize}$ and $\textsc{DfsHld}$ each visit every vertex and edge once: $\Oh(n)$ time,
$\Oh(n)$ space for the six arrays. Building the segment tree is $\Oh(n)$. Total preprocessing
$\Oh(n)$. (Recursion depth can be $\Th(n)$ on a path graph; in practice convert to an explicit stack
or raise the recursion limit.)
\end{complexitybox}

\begin{trapbox}
If your tree is given with a recursion depth that can reach $n$ (a caterpillar or a path), a recursive
$\textsc{DfsSize}$ can blow the stack at $n\approx 10^5$. Either iterate with an explicit stack, or
bump the system recursion limit. This is an implementation hazard, not an algorithmic one, but it
sinks more contest submissions than any logic bug in HLD.
\end{trapbox}

\section{Answering a path query}
\label{ch19:pathquery}

To aggregate over the path $u\!\to\!v$, repeatedly fold in the top of whichever endpoint sits in the
deeper chain, then jump to the chain top's parent. We measure ``deeper chain'' by the depth of the
chain \emph{top}: whichever of $\mathrm{top}[u],\mathrm{top}[v]$ is deeper gets processed, because its
chain cannot contain the other endpoint. We fold the contiguous segment from that top down to the
endpoint, hop above the top via its parent (a light edge), and continue until both endpoints share a
chain, where one final range from the shallower to the deeper finishes the job.

\begin{algorithm}[H]
\caption{\textsc{PathQuery}$(u,v)$ over vertex values}
\label{ch19:alg-query}
\DontPrintSemicolon
\KwIn{$u,v$; segment tree $S$ supporting range aggregation $\oplus$ over positions.}
\KwOut{$\bigoplus$ of values on the $u\!\to\!v$ path.}
$\mathrm{res}\gets \text{identity of }\oplus$\;
\While{$\mathrm{top}[u]\ne \mathrm{top}[v]$}{
  \If{$\mathrm{dep}[\mathrm{top}[u]] < \mathrm{dep}[\mathrm{top}[v]]$}{swap $u\leftrightarrow v$}
  \tcp{$\mathrm{top}[u]$ is the deeper chain top; fold its segment}
  $\mathrm{res}\gets \mathrm{res}\,\oplus\, S.\textsc{query}\big(\mathrm{pos}[\mathrm{top}[u]],\ \mathrm{pos}[u]\big)$\;
  $u\gets \mathrm{par}[\mathrm{top}[u]]$ \quad\tcp{climb the light edge above the chain}
}
\tcp{now $u,v$ share a chain; fold the segment between them}
\If{$\mathrm{dep}[u] > \mathrm{dep}[v]$}{swap $u\leftrightarrow v$}
$\mathrm{res}\gets \mathrm{res}\,\oplus\, S.\textsc{query}\big(\mathrm{pos}[u],\ \mathrm{pos}[v]\big)$\;
\Return $\mathrm{res}$\;
\end{algorithm}

\begin{whybox}
\textbf{Correctness.} Invariant: at the top of the loop, $\mathrm{res}$ holds the aggregate of all path
vertices already passed, and the still-unprocessed part of the path is exactly the path between the
\emph{current} $u$ and $v$. When $\mathrm{top}[u]\ne\mathrm{top}[v]$, the deeper top, say
$\mathrm{top}[u]$, cannot be an ancestor of $v$ (if it were, $v$ would lie in $u$'s chain or below a
shallower top), so the entire segment from $\mathrm{top}[u]$ down to $u$ belongs to the path and is a
contiguous index range $[\mathrm{pos}[\mathrm{top}[u]],\mathrm{pos}[u]]$. Folding it and moving $u$ to
$\mathrm{par}[\mathrm{top}[u]]$ removes exactly that segment from the remaining path and preserves the
invariant. The loop ends when both share a chain, i.e. both lie on one contiguous block; the final
range from the shallower to the deeper closes the path, whose shallower endpoint is by then the LCA. \qed
\end{whybox}

\begin{whybox}
\textbf{Why it is fast.} Each loop iteration moves one endpoint above its current chain top across a
\emph{light} edge. By Theorem~\ref{ch19:thm-loglight}, each of the two vertical halves
(endpoint up to LCA) crosses at most $\floor{\log_2 n}$ light edges, so the loop runs $\Oh(\log n)$
times. Each iteration issues one segment-tree range query at $\Oh(\log n)$. Total $\Oh(\log^2 n)$ per
path query.
\end{whybox}

\subsection{Tracing a query by hand}

\begin{labbox}
\textbf{Query the path $8 \to 7$ on the running tree} (Figure~\ref{ch19:fig-tree}), using the
positions and tops from Section~\ref{ch19:layout}. Recall
$\mathrm{top}[8]=8,\ \mathrm{top}[7]=7$, depths
$\mathrm{dep}[8]=3,\mathrm{dep}[7]=2$, and
$\mathrm{dep}[\mathrm{top}[8]]=3$, $\mathrm{dep}[\mathrm{top}[7]]=2$.
The true path is $8\!\to\!5\!\to\!2\!\to\!1\!\to\!3\!\to\!7$.

\emph{Iteration 1.} $\mathrm{top}[8]=8\ne 7=\mathrm{top}[7]$. Deeper top is $8$ (depth $3>2$). Fold
range $[\mathrm{pos}[8],\mathrm{pos}[8]]=[5,5]$, i.e. just vertex $8$. Climb:
$u\gets \mathrm{par}[\mathrm{top}[8]]=\mathrm{par}[8]=5$.

\emph{Iteration 2.} Now $u=5,\,v=7$: $\mathrm{top}[5]=1,\ \mathrm{top}[7]=7$, with
$\mathrm{dep}[\mathrm{top}[5]]=\mathrm{dep}[1]=0$, $\mathrm{dep}[\mathrm{top}[7]]=2$.
Deeper top is $7$, so swap to process the $v$-side: fold
$[\mathrm{pos}[7],\mathrm{pos}[7]]=[10,10]$ (vertex $7$). Climb:
$v\gets \mathrm{par}[7]=3$.

\emph{Iteration 3.} $u=5,\,v=3$: $\mathrm{top}[5]=1,\ \mathrm{top}[3]=3$,
depths of tops $0$ vs $1$. Deeper top is $3$; fold $[\mathrm{pos}[3],\mathrm{pos}[3]]=[8,8]$ (vertex
$3$). Climb: $v\gets \mathrm{par}[3]=1$.

\emph{Iteration 4.} $u=5,\,v=1$: $\mathrm{top}[5]=\mathrm{top}[1]=1$. Loop exits. Final segment:
$\mathrm{dep}[5]=2 > \mathrm{dep}[1]=0$, so swap; fold
$[\mathrm{pos}[1],\mathrm{pos}[5]]=[0,2]$, i.e. vertices $1,2,5$.

Collected ranges: $\{8\},\{7\},\{3\},\{1,2,5\}$, whose union is
$\{1,2,3,5,7,8\}$ -- exactly the path $8,5,2,1,3,7$ (Figure~\ref{ch19:fig-pathsplit}). Four range
queries, matching the
$\Oh(\log n)$ chain bound for $n=11$ ($\floor{\log_2 11}+1=4$). Each range query is a logarithmic
segment-tree call.
\end{labbox}

\begin{figure}[H]
\centering
\begin{tikzpicture}[
  every node/.style={circle,draw=nodeblue,thick,fill=blue!5,minimum size=8mm,inner sep=0pt,font=\small},
  onpath/.style={fill=structurepurple!25,draw=structurepurple,very thick},
  pathedge/.style={very thick,structurepurple},
  other/.style={edgegray,thin},
]
\node[onpath] (1) at (0,0) {1};
\node[onpath] (2) at (-3,-1.4) {2};
\node[onpath] (3) at (2.2,-1.4) {3};
\node (4) at (-4.2,-2.8) {4};
\node[onpath] (5) at (-1.8,-2.8) {5};
\node (6) at (1.0,-2.8) {6};
\node[onpath] (7) at (3.4,-2.8) {7};
\node[onpath] (8) at (-2.4,-4.2) {8};
\node (9) at (-1.2,-4.2) {9};
\node (10) at (0.4,-4.2) {10};
\node (11) at (-1.2,-5.6) {11};

\draw[pathedge] (1) -- (2);
\draw[pathedge] (1) -- (3);
\draw[other] (2) -- (4);
\draw[pathedge] (2) -- (5);
\draw[pathedge] (5) -- (8);
\draw[other] (5) -- (9);
\draw[other] (5) -- (10);
\draw[other] (9) -- (11);
\draw[other] (3) -- (6);
\draw[pathedge] (3) -- (7);
\end{tikzpicture}
\caption{The path $8\!\to\!7$ (purple vertices) split into the four folded ranges $\{8\}$, $\{1,2,5\}$,
$\{3\}$, $\{7\}$. The block $\{1,2,5\}$ is a stretch of chain A; the singletons $\{8\},\{3\},\{7\}$ are
the bottoms of their own chains. The light hops $8\!\to\!5$ and $7\!\to\!3\!\to\!1$ are where the loop
jumps from one chain to the next.}
\label{ch19:fig-pathsplit}
\end{figure}

\section{Path updates and lazy propagation}
\label{ch19:pathupdate}

A path \emph{update} is identical in structure: decompose the path into the same $\Oh(\log n)$
contiguous ranges and apply a range update to each, using a lazy-propagating segment tree. Replace
each $S.\textsc{query}(\cdot,\cdot)$ in Algorithm~\ref{ch19:alg-query} with
$S.\textsc{rangeUpdate}(\cdot,\cdot,\delta)$ and drop the accumulator.

\begin{algorithm}[H]
\caption{\textsc{PathUpdate}$(u,v,\delta)$ -- add $\delta$ on the path}
\label{ch19:alg-update}
\DontPrintSemicolon
\While{$\mathrm{top}[u]\ne \mathrm{top}[v]$}{
  \If{$\mathrm{dep}[\mathrm{top}[u]] < \mathrm{dep}[\mathrm{top}[v]]$}{swap $u\leftrightarrow v$}
  $S.\textsc{rangeUpdate}\big(\mathrm{pos}[\mathrm{top}[u]],\ \mathrm{pos}[u],\ \delta\big)$\;
  $u\gets \mathrm{par}[\mathrm{top}[u]]$\;
}
\If{$\mathrm{dep}[u] > \mathrm{dep}[v]$}{swap $u\leftrightarrow v$}
$S.\textsc{rangeUpdate}\big(\mathrm{pos}[u],\ \mathrm{pos}[v],\ \delta\big)$\;
\end{algorithm}

\begin{complexitybox}
$\Oh(\log n)$ chain segments, each a lazy range update at $\Oh(\log n)$: $\Oh(\log^2 n)$ per path
update. With a Fenwick tree and a difference trick you can sometimes shave the inner factor for
specific aggregates, but $\Oh(\log^2 n)$ is the standard, safe bound.
\end{complexitybox}

\subsection{Edge values vs vertex values}

Many problems put values on \emph{edges}, not vertices. Map each non-root vertex $v$ to its parent
edge $(\mathrm{par}[v],v)$ and store that edge's value at $\mathrm{pos}[v]$. The root owns no parent
edge, so its slot stays at the identity. The only change to a path query/update is the final
within-chain range: after the loop, $u$ is the LCA, and the LCA's vertex slot is \emph{not} an edge of
the path, so fold $[\mathrm{pos}[u]+1,\ \mathrm{pos}[v]]$ instead of $[\mathrm{pos}[u],\mathrm{pos}[v]]$
(and skip entirely if $u=v$).

\begin{trapbox}
The off-by-one at the LCA is the single most common HLD bug. For \emph{vertex} values, the LCA is on
the path and its slot must be included. For \emph{edge} values, the LCA's slot represents the edge
\emph{above} the LCA, which is off the path, so it must be excluded with the $+1$. Mixing these up
gives answers that are wrong by exactly one vertex/edge and are hard to spot because short paths
look correct.
\end{trapbox}

\section{Subtree queries and updates}
\label{ch19:subtree}

Subtrees are easier than paths, because of the contiguity we already proved: the subtree of $v$ is the
position range $[\mathrm{pos}[v],\ \mathrm{pos}[v]+\mathrm{sz}[v]-1]$. So a subtree query or update is a
\emph{single} range operation -- no chain walking at all.

\begin{whybox}
The heavy-first DFS is still a preorder, so when it stamps $v$ at position $\mathrm{pos}[v]$ it then
stamps every descendant of $v$ before returning, giving them positions
$\mathrm{pos}[v]+1,\dots,\mathrm{pos}[v]+\mathrm{sz}[v]-1$ contiguously. Hence
$\textsc{SubtreeQuery}(v)=S.\textsc{query}(\mathrm{pos}[v],\ \mathrm{pos}[v]+\mathrm{sz}[v]-1)$ and the
analogous range update. Both cost $\Oh(\log n)$.
\end{whybox}

\begin{labbox}
On the running tree, subtree of $5$ has $\mathrm{pos}[5]=2,\ \mathrm{sz}[5]=5$, so range
$[2,6]=\{5,9,11,8,10\}$ -- correct. Subtree of $2$: $\mathrm{pos}[2]=1,\ \mathrm{sz}[2]=7$, range
$[1,7]=\{2,5,9,11,8,10,4\}$. A single segment-tree call covers either, no per-chain loop.
\end{labbox}

\begin{aestheticbox}
One array layout, two payoffs. The same $\mathrm{pos}$ ordering that makes chains contiguous (for
$\Oh(\log^2 n)$ path operations) also makes subtrees contiguous (for $\Oh(\log n)$ subtree
operations). That is not a coincidence: heavy-first DFS is a preorder, and \emph{every} preorder makes
subtrees contiguous; heavy-first is the special preorder that \emph{additionally} aligns chains. You
pay nothing extra for the subtree power.
\end{aestheticbox}

\section{Complexity, all in one place}
\label{ch19:complexity}

\begin{complexitybox}
Let $n$ be the number of vertices and $Q$ the number of operations.
\begin{center}
\begin{tabular}{lll}
\toprule
Operation & Time & Why \\
\midrule
Preprocessing (two DFS + segtree build) & $\Oh(n)$ & each vertex/edge touched $\Oh(1)$ times\\
Memory & $\Oh(n)$ & six arrays + segment tree\\
Path query / update & $\Oh(\log^2 n)$ & $\Oh(\log n)$ chains $\times$ $\Oh(\log n)$ segtree\\
Subtree query / update & $\Oh(\log n)$ & one contiguous range\\
Point query / update & $\Oh(\log n)$ & one segtree leaf path\\
LCA via the path loop & $\Oh(\log n)$ & last shared top is the LCA\\
\bottomrule
\end{tabular}
\end{center}
Total over $Q$ operations: $\Oh(n + Q\log^2 n)$. The two $\log$ factors come from two independent
logarithmic structures stacked: the HLD chain count and the segment tree depth.
\end{complexitybox}

\begin{whybox}
\textbf{Can we drop a log?} For commutative-group aggregates (sum), pairing HLD with a Fenwick tree
keeps $\Oh(\log^2 n)$ but with smaller constants. For \emph{static} path \emph{queries} (no updates),
replace the segment tree by an $\Oh(1)$ sparse table over each chain and the cost drops to
$\Oh(\log n)$ per query -- $\Oh(\log n)$ chains, $\Oh(1)$ each. With updates the inner $\log$ is
genuinely needed, so $\Oh(\log^2 n)$ is the right target for the dynamic case. A different
decomposition, link-cut trees (Chapter~22), achieves $\Oh(\log n)$ \emph{amortized} per path operation
even with structural changes, at the cost of much heavier machinery.
\end{whybox}

\section{Finding the LCA for free}
\label{ch19:lca}

The path loop already walks both endpoints up to their meeting chain, so it computes
$\mathrm{LCA}(u,v)$ as a side effect: when the loop exits, the shallower of $u,v$ is the LCA.

\begin{algorithm}[H]
\caption{\textsc{Lca}$(u,v)$ via heavy paths}
\label{ch19:alg-lca}
\DontPrintSemicolon
\While{$\mathrm{top}[u]\ne \mathrm{top}[v]$}{
  \If{$\mathrm{dep}[\mathrm{top}[u]] < \mathrm{dep}[\mathrm{top}[v]]$}{swap $u\leftrightarrow v$}
  $u\gets \mathrm{par}[\mathrm{top}[u]]$\;
}
\lIf{$\mathrm{dep}[u]\le\mathrm{dep}[v]$}{\Return $u$} \lElse{\Return $v$}
\end{algorithm}

\begin{labbox}
$\textsc{Lca}(8,7)$: tops $8,7$, deeper $8\Rightarrow u\gets\mathrm{par}[8]=5$. Now tops $1,7$, deeper
$7\Rightarrow v\gets\mathrm{par}[7]=3$. Tops $1,3$, deeper $3\Rightarrow v\gets\mathrm{par}[3]=1$. Tops
$1,1$ equal, exit. Shallower of $5,1$ is $1$. Indeed $\mathrm{LCA}(8,7)=1$, matching the trace in
Section~\ref{ch19:pathquery}. Cost $\Oh(\log n)$: pure pointer climbing, no segment tree.
\end{labbox}

\section{Worked end-to-end example with values}
\label{ch19:fullexample}

Put a value on each vertex and run a real mix of operations, so the segment tree side is concrete.
Initialize vertex $i$'s value to $i$ (so $\mathrm{val}=[1,\dots,11]$ indexed by vertex). Lay them by
position using Section~\ref{ch19:layout}'s table; the segment-tree array $a[\,]$ indexed by position is

\begin{center}
\begin{tabular}{c|ccccccccccc}
position $p$ & 0 & 1 & 2 & 3 & 4 & 5 & 6 & 7 & 8 & 9 & 10\\\hline
vertex      & 1 & 2 & 5 & 9 & 11& 8 & 10& 4 & 3 & 6 & 7\\
$a[p]$      & 1 & 2 & 5 & 9 & 11& 8 & 10& 4 & 3 & 6 & 7\\
\end{tabular}
\end{center}

\begin{labbox}
\textbf{(1) Sum on path $11\to 10$.} Path is $11\!\to\!9\!\to\!5\!\to\!10$. Run the loop. Initially
$\mathrm{top}[11]=1$ (depth $0$), $\mathrm{top}[10]=10$ (depth $3$): deeper top $10$, swap so we fold
the $10$-side: range $[\mathrm{pos}[10],\mathrm{pos}[10]]=[6,6]$, value $a[6]=10$; climb
$v\gets\mathrm{par}[10]=5$. Now $u=11,v=5$, tops $1$ and $1$ equal -- exit. Final range:
$\mathrm{dep}[11]=4>\mathrm{dep}[5]=2$ so swap; fold $[\mathrm{pos}[5],\mathrm{pos}[11]]=[2,4]$,
values $a[2]+a[3]+a[4]=5+9+11=25$. Total $=10+25=35$. Check directly:
$\mathrm{val}(11)+\mathrm{val}(9)+\mathrm{val}(5)+\mathrm{val}(10)=11+9+5+10=35.\ \checkmark$

\textbf{(2) Add $\delta=100$ on path $4\to 6$.} Path is $4\!\to\!2\!\to\!1\!\to\!3\!\to\!6$. Loop:
tops $\mathrm{top}[4]=4$ (depth $2$), $\mathrm{top}[6]=3$ (depth $1$): deeper top $4$, range
$[\mathrm{pos}[4],\mathrm{pos}[4]]=[7,7]$ gets $+100$ (vertex $4$); climb $u\gets\mathrm{par}[4]=2$.
Now $u=2$ ($\mathrm{top}=1$, depth $0$), $v=6$ ($\mathrm{top}=3$, depth $1$): deeper top $3$, swap;
range $[\mathrm{pos}[3],\mathrm{pos}[6]]=[8,9]$ gets $+100$ (vertices $3,6$); climb
$v\gets\mathrm{par}[3]=1$. Now $u=2,v=1$ tops $1,1$ equal -- exit. Final: $\mathrm{dep}[2]=1>0$, swap;
range $[\mathrm{pos}[1],\mathrm{pos}[2]]=[0,1]$ gets $+100$ (vertices $1,2$). Updated vertices:
$\{4,3,6,1,2\}$, exactly the path. New values: $a[0]=101,a[1]=102,a[7]=104,a[8]=103,a[9]=106$.

\textbf{(3) Sum on subtree of $2$.} Range $[\mathrm{pos}[2],\mathrm{pos}[2]+\mathrm{sz}[2]-1]=[1,7]$.
After update (2) this is $a[1{:}7]=102+5+9+11+8+10+104=249$. Vertices $\{2,5,9,11,8,10,4\}$; original
sum $2+5+9+11+8+10+4=49$, and the two of them on path $4\to6$ that got $+100$ are $2$ and $4$, i.e.
$+200$, giving $49+200=249.\ \checkmark$ (Vertices $1,3,6$ are not in the subtree of $2$.)
\end{labbox}

\section{Reference implementation sketch}
\label{ch19:impl}

The pieces fit into a compact program: two DFS for the layout, a lazy segment tree for the array, and
the chain-walking wrappers. The skeleton below assumes a $0$-indexed segment tree exposing
$\textsc{query}(l,r)$ and $\textsc{update}(l,r,\delta)$ over inclusive ranges.

\begin{lstlisting}[language=C++]
// par[], dep[], sz[], heavy[], top[], pos[] filled by Algorithm 19.1.
long long pathQuery(int u, int v) {
    long long res = 0;                      // identity of the aggregate
    while (top[u] != top[v]) {
        if (dep[top[u]] < dep[top[v]]) swap(u, v);
        res += seg.query(pos[top[u]], pos[u]);
        u = par[top[u]];
    }
    if (dep[u] > dep[v]) swap(u, v);
    res += seg.query(pos[u], pos[v]);       // pos[u]+1 for EDGE values
    return res;
}
void pathUpdate(int u, int v, long long d) {
    while (top[u] != top[v]) {
        if (dep[top[u]] < dep[top[v]]) swap(u, v);
        seg.update(pos[top[u]], pos[u], d);
        u = par[top[u]];
    }
    if (dep[u] > dep[v]) swap(u, v);
    seg.update(pos[u], pos[v], d);          // pos[u]+1 for EDGE values
}
long long subtreeQuery(int v) {             // contiguous range, one call
    return seg.query(pos[v], pos[v] + sz[v] - 1);
}
\end{lstlisting}

\begin{trapbox}
Three recurring pitfalls. (i) Comparing $\mathrm{dep}[u]$ instead of $\mathrm{dep}[\mathrm{top}[u]]$
inside the loop -- the loop condition is about which \emph{chain} is deeper, decided by the top's
depth. (ii) The edge-value $+1$ at the LCA, already flagged. (iii) Forgetting that the final range
needs the shallower endpoint first; if $\mathrm{pos}[u]>\mathrm{pos}[v]$ your segment tree may silently
return the identity or crash. Swap by depth before the final call.
\end{trapbox}

\section{Where it gets used}
\label{ch19:apps}

\begin{applicationbox}
\textbf{Network paths.} ``Maximum congested link on the route between two routers, with link costs
changing over time'' is a path-max query with point updates -- HLD plus a max segment tree.

\textbf{Accumulation on org charts and phylogenies.} ``Add a bonus to everyone on the chain of command
from employee $u$ up to manager $v$'' is a path add; ``total payroll under manager $w$'' is a subtree
sum. Both are exactly the operations above.

\textbf{Competitive programming.} HLD is the standard answer to tree path queries with updates
(e.g.\ SPOJ \texttt{QTREE}, many Codeforces tree problems). When the values are on edges and the query
is path-max, HLD over a max segment tree is the textbook solution.

\textbf{Auxiliary structure.} HLD underlies fast implementations of tree-path versions of problems
that on arrays use segment trees: range assignment, range gcd, range affine transforms -- each lifts
to trees by swapping the underlying segment tree.
\end{applicationbox}

\begin{historybox}
The size-balanced heavy-path idea traces to Sleator and Tarjan's 1983 paper on \emph{link-cut trees},
where heavy paths (there called the \emph{preferred paths}) are restructured dynamically to give
$\Oh(\log n)$ amortized path operations even as the tree changes shape. Harel and Tarjan's 1984
LCA work used a related path partition. ``Heavy-light decomposition'' as a static, segment-tree-backed
technique became a competitive-programming staple in the 2000s precisely because it reduces the hard
part (tree paths) to the easy part (array ranges) you already know. The $\floor{\log_2 n}$ light-edge
bound is the same counting argument that powers small-to-large merging and many divide-and-conquer
tree algorithms.
\end{historybox}

\begin{connectionbox}
\textbf{Range structures (Chapter~23).} HLD is a reduction: it converts tree-path problems into the
array-range problems solved by segment trees, Fenwick trees, and sparse tables. Everything you know
about lazy propagation, range-assign, and merge-sort trees transfers verbatim to tree paths.

\textbf{Euler tour technique (Chapter~21).} The heavy-first DFS \emph{is} an Euler/preorder tour; the
subtree-contiguity we used is the same in/out-time trick that the Euler tour chapter develops for
subtree queries. HLD adds the chain alignment on top.

\textbf{Link-cut trees (Chapter~22).} LCTs maintain heavy (preferred) paths dynamically under
\texttt{link}/\texttt{cut}, achieving $\Oh(\log n)$ amortized where static HLD needs the tree fixed.
Static HLD is the right tool when the tree shape never changes; LCTs when it does.

\textbf{Centroid decomposition (Chapter~20).} A sibling reduction that also yields $\Oh(\log n)$
recursion depth, but partitions by \emph{centroids} for path-\emph{counting} problems rather than
laying paths into an array. The two solve different query shapes; knowing both lets you pick.
\end{connectionbox}

\section{Exercises}
\label{ch19:exercises}

\begin{exercise}[Trace another path]
On the running tree (Figure~\ref{ch19:fig-tree}), trace \textsc{PathQuery}$(4,11)$ step by step. List
every chain range folded, the value of $u,v$ after each iteration, and confirm the union of ranges is
the true path $4\!\to\!2\!\to\!5\!\to\!9\!\to\!11$.
\end{exercise}

\begin{exercise}[Edge values]
Reimplement \textsc{PathQuery} for values stored on \emph{edges} (each edge's value at the child's
position). Prove that the only change is the final within-chain range, which becomes
$[\mathrm{pos}[u]+1,\mathrm{pos}[v]]$, and that the result is empty when $u=v$. Trace path-max on edges
for the pair $(8,7)$ assuming each edge $(p,c)$ has weight $\mathrm{dep}[c]$.
\end{exercise}

\begin{exercise}[Break the tie rule]
Show that choosing the heavy child by \emph{smallest} subtree instead of largest can violate
Lemma~\ref{ch19:lem-halve}. Give a concrete small tree where the light-edge count on some root-to-leaf
path exceeds $\floor{\log_2 n}$ under the wrong rule, and explain which inequality in the proof fails.
\end{exercise}

\begin{exercise}[Subtree add, path sum]
Design a structure supporting subtree-add updates and path-sum queries together. Argue that a single
lazy segment tree over the HLD positions handles both, since a subtree is one range and a path is
$\Oh(\log n)$ ranges. State the per-operation cost and prove correctness of combining the two range
families on one tree.
\end{exercise}

\begin{exercise}[Static queries, one log]
Suppose there are no updates, only path-max queries. Replace the segment tree with a sparse table over
the whole position array. Prove that each path query then costs $\Oh(\log n)$ total, not
$\Oh(\log^2 n)$, and identify exactly where the second log disappears.
\end{exercise}

\begin{exercise}[Kth vertex on a path]
Extend HLD to answer ``the $k$-th vertex on the path from $u$ to $v$'' in $\Oh(\log n)$. Hint: split at
the LCA using depths, then locate within chains by position arithmetic. Prove your index computation
handles the LCA boundary correctly.
\end{exercise}

\begin{exercise}[Distance and the LCA loop]
Using only $\mathrm{dep}$ and Algorithm~\ref{ch19:alg-lca}, give an $\Oh(\log n)$ formula for the
number of edges on the path $u\!\to\!v$. Prove it from
$\dist(u,v)=\mathrm{dep}[u]+\mathrm{dep}[v]-2\,\mathrm{dep}[\mathrm{LCA}(u,v)]$.
\end{exercise}

\begin{exercise}[Worst case is one chain]
Prove that on a path graph (a line of $n$ vertices) the entire tree is a single heavy chain, so every
path query is a \emph{single} segment-tree range and costs $\Oh(\log n)$, not $\Oh(\log^2 n)$.
Conversely, construct a tree where some path query truly forces $\Th(\log n)$ distinct chain ranges,
showing the $\Oh(\log^2 n)$ bound is tight in the worst case.
\end{exercise}

\begin{exercise}[Iterative build]
Rewrite \textsc{DfsSize} and \textsc{DfsHld} (Algorithm~\ref{ch19:alg-build}) iteratively with an
explicit stack so they survive a path graph of $n=10^6$ without stack overflow. Verify your positions
and tops match the recursive version on the running tree.
\end{exercise}

\begin{exercise}[Forest, shared array]
You are given a forest. Adapt the layout so all components share one position array (one segment
tree), with each component's vertices contiguous. Show path queries within a component still cost
$\Oh(\log^2 n)$ and that no query ever crosses components.
\end{exercise}
